% EGA V, English translation by P. Blass, 1987 (cote 24): archive cote 24, site page 33 (PDF page 34).
% typescript. Draft compared: none.
% AI-transcribed by Claude Opus 5.5 (claude-opus-5-5, adaptive thinking, effort high), Message Batches API, batch ega24b (msgbatch_01RGJJeHdR96U16ZWFuUjtDx), 2026-11-13; not yet checked by eye.
\scanpage{24/33}
\mspage{33}

and if we denote by E the set of $\xi \in \overset{\vee}{\underline{P}}$ such that $G_\xi$ resp. $\underline{Y}_\xi$
satisfies $\underline{P}$ except over a closed set of dimension $\leq$ m (i.e. the
set $\underline{P}$-singular Z is of dimension $\leq$ m).  Then a) E is a
constructible subset of $\underline{P}\mathrm{v}$ and b) in the case S = Spec(k) if F
respectively X satisfies $\underline{P}$ except over a set of dimension $\leq$ m + 1
then E contains a non-empty open set.

\uline{Proposition 7.2}

Let us assume that X is proper over S and that F is of finite
presentation finite and flat over S.  Let $\underline{P}$ be one of the
properties (i) to (v) of (4.2) and let m be an integer.  Let s $\in$ S
and let us suppose that the set $Z_s$ of points of $X_s$ where $F_s$ does
not satisfy $\underline{P}$ is of dimension $\leq$ m + 1.  Then if also k(s) is
infinite there exists a neighbourhood U of s in E and a section $\xi$
of $\underline{P}\mathrm{v}$ over U having the following properties :  For every s$\in$ U the
set of points of $Y_{\xi,\ s}$ where $G_{\xi(s)}$ does not satisfy $\underline{P}$ is of
dimension $\leq$ m and $\phi_{\xi(s)}$ is $F_{\xi,s}$ regular.  Under such conditions the
module $G_\xi$ over $Y_\xi$ is flat relative to U.  Finally if k(s) is not
supposed infinite we can again make the previous construction
after an étale extension of the type anticipated in 7.1.

\uline{Proposition 7.3}.  Essentially the same.  There is no longer an F
and assume that X is flat relative to S we refer to properties (i)
to (v) of 4.3 in place of those of 4.2 (but being careful to make
the reservation.")  k(s) of characteristic zero or f: X $\rightarrow$ $\underline{P}$ is an
